\begin{document}

% ------------------------------------------------------------
% 表紙
% ------------------------------------------------------------
\begin{titlepage}
\begin{tikzpicture}[remember picture,overlay]
  \fill[bkblue] (current page.north west) rectangle ([yshift=-92mm]current page.north east);
  \fill[bkorange] ([yshift=-92mm]current page.north west) rectangle ([yshift=-95mm]current page.north east);
\end{tikzpicture}
\vspace*{8mm}
\begin{center}
{\color{white}\large 統計検定1級　統計応用（医薬生物学）対策\par}
\vspace{10mm}
{\color{white}\fontsize{30}{40}\selectfont\bfseries 医療統計学参考書\par}
\vspace{6mm}
{\color{white}\large データ構造から手法を選び，導出し，解釈する\par}
\end{center}
\vspace{40mm}
\begin{center}
\begin{tikzpicture}[node distance=4mm,
  b/.style={draw=bkblue,rounded corners=2pt,fill=bkblue!6,minimum width=22mm,minimum height=10mm,align=center,font=\small}]
  \node[b](a){データ構造};
  \node[b,right=of a](b){推定対象};
  \node[b,right=of b](c){手法};
  \node[b,right=of c](d){仮定};
  \node[b,right=of d](e){計算};
  \node[b,right=of e](f){解釈};
  \foreach \x/\y in {a/b,b/c,c/d,d/e,e/f}\draw[-{Stealth[length=2.5mm]},thick,bkblue](\x)--(\y);
\end{tikzpicture}
\end{center}
\vspace{30mm}
\begin{center}
\begin{tabular}{l}
過去問9回分（2016--2019，2021--2025）の分析に基づく\\
生存時間解析・分割表・診断・交絡・臨床試験デザイン・メタアナリシス・欠測・ベイズ
\end{tabular}
\end{center}
\vfill
\begin{center}
{\small 学習用電子参考書　2026年版}
\end{center}
\end{titlepage}

\frontmatter
\pagenumbering{roman}

% ------------------------------------------------------------
\chapter*{はじめに}
\addcontentsline{toc}{chapter}{はじめに}
\markboth{はじめに}{はじめに}

本書は，統計検定1級「統計応用（医薬生物学）」の受験者のための参考書である．
医療統計の入門書は多いが，1級の記述式問題は「手法の名前を知っている」だけでは解けない．
典型的な設問は次の4段で構成されている．
\begin{enumerate}[label=(\arabic*)]
\item 問題文の研究デザインとデータ構造を読み取り，
\item 推定したい量（効果指標）と，それに合う検定・推定の方法を選び，
\item 尤度・分散・漸近分布などを導出して数値を計算し，
\item 結果を医学・疫学・臨床試験の文脈で解釈する．
\end{enumerate}
本書はこの4段をそのまま章の骨格にしている．
各テーマを
\[
\boxed{\ \text{データ構造}\ \rightarrow\ \text{推定対象}\ \rightarrow\ \text{手法}\ \rightarrow\ \text{仮定}\ \rightarrow\ \text{計算}\ \rightarrow\ \text{解釈}\ }
\]
の流れで説明する．

内容は2016年から2025年までの9回分（2020年は実施なし）の医薬生物学の過去問と，
同じ年度の統計数理の過去問を分析して設計した．
設問ごとに主テーマ・副テーマ・必要能力を分類し，出題の厚い分野に紙幅を多く割いている．
ただし「過去に出ていないから不要」とは考えず，標準的な医療統計として最低限必要な内容は補った．
過去問で実際に問われた内容と，一般的に重要だが直接は問われていない内容は，
「過去問での出題」ボックスと「補足」「発展」の表示で区別している．

過去問の本文は転載せず，「\kakomon{2022}{2}」のように年と問番号で参照する．
例題はすべて，過去問と同じ能力を問うように作ったオリジナル問題である．

% ------------------------------------------------------------
\section*{本書の対象読者}
\addcontentsline{toc}{section}{本書の対象読者}
次の読者を想定している．
\begin{itemize}
\item 統計検定準1級程度の内容を習得している．
\item 『現代数理統計学の基礎』（久保川達也）第2〜7章，すなわち確率分布，多次元分布，標本分布と極限定理，推定，検定を学習済みである．
\item 最尤法，Fisher情報量，デルタ法，漸近正規性，尤度比検定，基本的な区間推定・検定を，ある程度自分で使える．
\end{itemize}
したがって「期待値とは」「最尤法とは」といった一般の数理統計は一から説明しない．
医療統計の議論で必要になった箇所だけを「数理統計との接続」として補う．
逆に，順位統計量，打ち切りのある尤度，部分尤度，多変量デルタ法，分割表固有の方法，
多重性の理論などは上記の範囲に含まれないので，本書の中で導出する．

% ------------------------------------------------------------
\section*{本書の使い方}
\addcontentsline{toc}{section}{本書の使い方}
\paragraph{章の構成}
各章は原則として次の順に並ぶ．内容に応じて一部を省略・統合している．
\begin{center}\small
\begin{tabular}{@{}ll@{}}
\toprule
1. この章でできるようになること & 8. 結果の解釈\\
2. 過去問での出題 & 9. 手法選択\\
3. 問題設定 & 10. よくある誤り\\
4. 基本概念 & 11. 数理統計との接続\\
5. 数理 & 12. オリジナル例題\\
6. 推定・検定 & 13. 解答\\
7. 仮定 & 14. 章末まとめ\\
\bottomrule
\end{tabular}
\end{center}

\paragraph{ボックスと表示}
\begin{center}\small
\begin{tabular}{@{}lL{110mm}@{}}
\toprule
表示 & 意味\\
\midrule
\Hissu & 試験で導出・計算できなければならない内容\\
\Hinshutsu & 9回の中で複数回出題された内容\\
\Youchui & 取り違えやすい概念・符号・自由度など\\
\Hatten & 試験に直接出にくいが理解を深める内容\\
\textcolor{bkorange}{\textbf{過去問での出題}} & その章のテーマが実際にどの年・問で問われたか\\
\textcolor{bkred}{\textbf{よくある誤り}} & 答案で起こりやすい誤り\\
\textcolor{bkpurple}{\textbf{数理統計との接続}} & 『現代数理統計学の基礎』や統計数理の過去問との対応\\
\textcolor{bkgray}{\textbf{補足・発展}} & 過去問・指定教材に直接ない内容（Web資料等から補った場合は出典を付記）\\
\bottomrule
\end{tabular}
\end{center}

\paragraph{記法}
本書全体で次の記法を用いる．
\begin{itemize}
\item $z_\alpha$ は標準正規分布の上側 $100\alpha\%$ 点（$\Prob(Z>z_\alpha)=\alpha$）．
      $t_\nu(\alpha)$，$\chi^2_\nu(\alpha)$ も上側点．よく使う値は $z_{0.05}=1.645$，$z_{0.025}=1.960$，$z_{0.10}=1.282$，$z_{0.20}=0.842$．
\item $\beta$ は第2種の過誤確率，検出力は $1-\beta$．
      \Youchui 『現代数理統計学の基礎』では $\beta(\theta)$ を検出力関数そのものとして用いるので逆になる．
\item 不偏分散 $s^2=\frac{1}{n-1}\sum(x_i-\bar x)^2$．
      \Youchui 同書では $S^2$ が除数 $n$，$V^2$ が除数 $n-1$ である．
\item 指数分布 $\Ex(\lambda)$ は率（ハザード）$\lambda$ による表記で，密度 $\lambda e^{-\lambda t}$，平均 $1/\lambda$．
\item ガンマ分布 $\Ga(a,b)$ は\textbf{形状 $a$・率 $b$}（密度 $\frac{b^a}{\Gamma(a)}x^{a-1}e^{-bx}$，平均 $a/b$，分散 $a/b^2$）．
      試験（\kakomon{2025}{2}，\kakomonSM{2022}{3}）と同じ流儀である．
      \Youchui 『現代数理統計学の基礎』は $\Ga(\alpha,\beta)$ を\textbf{形状・尺度}（平均 $\alpha\beta$）で定義している．
\item $\Normal(\mu,\sigma^2)$ の第2引数は分散．$\log$ は自然対数．
\item $I_n(\theta)$ は標本全体のFisher情報量，$I_1(\theta)$ は1観測あたり．
\end{itemize}

\enlargethispage{2\baselineskip}
\paragraph{数値計算について}
統計検定1級では関数電卓が使えない．本書の例題の数値計算は，四則演算・平方根と巻末付表（正規分布表・$t$ 分布表・$\chi^2$ 分布表・指数関数と対数の表）で実行できる形で示し，
付表を引いた箇所では「$t_9(0.05)=1.833$」のようにどの値を使ったかを明記する．

% ------------------------------------------------------------
\chapter*{統計検定1級・医薬生物学の出題構造}
\addcontentsline{toc}{chapter}{統計検定1級・医薬生物学の出題構造}
\markboth{出題構造}{出題構造}

\paragraph{試験の形式}
統計検定1級は「統計数理」と「統計応用」の2科目からなり，それぞれ90分の記述式試験である．
統計応用は人文科学・社会科学・理工学・医薬生物学の4分野から1つを選んで受験し，
大問5問のうち3問を選択して解答する（1問あたり約30分）\cite{web:lec,web:academaid}．
医薬生物学の大問のうち問1〜問4が分野固有の問題，問5は4分野共通の問題である
（本書で分析した9回すべてでこの構成であった）．

\paragraph{9回分の傾向}
医薬生物学の問1〜4（計36問）を主テーマで分類すると，おおよそ次のようになる
（1問が複数テーマにまたがる場合は重複して数える）．
\begin{center}\small
\begin{tabular}{@{}L{62mm}cL{80mm}@{}}
\toprule
分野 & 出題数 & 代表的な問\\
\midrule
生存時間解析（KM・ログランク・指数モデル・Cox・RMST） & 7 & \kakomon{2016}{3}，\kakomon{2019}{1}，\kakomon{2022}{1}\\
2 値データ・分割表・McNemar・MH & 7 & \kakomon{2018}{2}，\kakomon{2021}{3}，\kakomon{2022}{2}\\
診断検査・ROC & 5 & \kakomon{2018}{3}，\kakomon{2019}{3}，\kakomon{2024}{1}\\
交絡・層別・標準化・傾向スコア & 4 & \kakomon{2017}{4}，\kakomon{2019}{2}\\
中間解析・二段階・多重性・用量反応 & 5 & \kakomon{2017}{2}，\kakomon{2024}{4}，\kakomon{2025}{4}\\
縦断データ・混合モデル・ベースライン調整 & 3 & \kakomon{2022}{3}，\kakomon{2023}{4}，\kakomon{2024}{2}\\
ロジスティック回帰 & 3 & \kakomon{2018}{4}，\kakomon{2023}{1}\\
ノンパラメトリック（順位）検定 & 3 & \kakomon{2016}{1}，\kakomon{2019}{4}，\kakomon{2025}{1}\\
検出力・症例数・イベント数 & 3 & \kakomon{2018}{1}，\kakomon{2023}{2}\\
同等性・生物学的同等性 & 2 & \kakomon{2016}{4}，\kakomon{2021}{2}\\
競合リスク & 2 & \kakomon{2021}{1}，\kakomon{2025}{3}\\
メタアナリシス & 2 & \kakomon{2021}{4}，\kakomon{2023}{3}\\
ベイズ・ポアソン率 & 2 & \kakomon{2022}{4}，\kakomon{2025}{2}\\
欠測データ・多重補完 & 1 & \kakomon{2024}{3}\\
\bottomrule
\end{tabular}
\end{center}
生存時間解析は9回中7回で何らかの形で出題されており，最重要分野である．
次いで2 値データ（分割表・診断・ロジスティック）が厚い．
2021年以降は，競合リスク，メタアナリシスの変量効果，混合モデル，多重補完，多重性，ベイズ予測など，
臨床試験の方法論に寄った出題が増えている．

\enlargethispage{4\baselineskip}
\paragraph{問われる能力}
設問のほとんどは「定義を式で書く→導出する→データで計算する→解釈を述べる」の組み合わせである．
導出の多くは，尤度の最大化，分散の計算（多項分布の共分散・条件付き分散を含む），
デルタ法，正規近似から検出力式を立てる，といった数理統計の基本操作に帰着する．
統計数理の過去問は医薬固有のモデル（打ち切り，Cox，分割表）を直接は扱わないが，
指数分布・ポアソン分布・ガンマ混合・デルタ法・Fisher情報などを
医薬の問題より深く問う（詳細は各章の「数理統計との接続」）．

% ------------------------------------------------------------
\chapter*{過去問9回分の出題マップ}
\addcontentsline{toc}{chapter}{過去問9回分の出題マップ}
\markboth{出題マップ}{出題マップ}

表中の数字は医薬生物学の問番号，「5c」は4分野共通の問5で内容が関連するものを表す．
設問単位の詳細な分類（副テーマ・必要能力）は\cref{app:D}の過去問対応表にある．

\begin{center}\small
\setlength{\tabcolsep}{3.2pt}
\begin{tabular}{@{}clccccccccc@{}}
\toprule
章 & テーマ & 2016 & 2017 & 2018 & 2019 & 2021 & 2022 & 2023 & 2024 & 2025\\
\midrule
\ref{ch:01} & 問題の読み方・P 値・CI & 2,5c & & & & & 1 & & & \\
\ref{ch:02} & 研究デザイン・効果指標 & & 4,5c & & 2 & & 2 & & & \\
\ref{ch:03} & 連続・ノンパラメトリック & 1,5c & & & 4 & & & & & 1\\
\ref{ch:04} & 2 値データ・分割表 & 2 & 4 & 2 & 3 & 3 & 2 & & 1 & \\
\ref{ch:05} & 診断検査・ROC & & & 3,4 & 3 & 5c & & 1 & 1 & \\
\ref{ch:06} & ロジスティック回帰 & & & 4 & 2 & & & 1 & & \\
\ref{ch:07} & 交絡・層別・マッチング & & 4 & & 2 & 3 & 2 & & & \\
\ref{ch:08} & 縦断・ベースライン・混合 & & & & & & 3,5c & 4 & 2 & \\
\ref{ch:09} & 生存時間解析の基礎 & 3 & 1 & & 1 & 1 & 1 & & & \\
\ref{ch:10} & Cox・RMST・パラメトリック & 3 & 1 & 1 & 1 & & 1 & 2 & & \\
\ref{ch:11} & 競合リスク & & & & & 1 & & & & 3\\
\ref{ch:12} & 検出力・症例数 & & 5c & 1 & & & & 2,5c & & 1\\
\ref{ch:13} & 同等性・生物学的同等性 & 4 & & & & 2 & & & & \\
\ref{ch:14} & 中間解析・多重性・用量反応 & 2 & 2,3 & & & & & & 4 & 4\\
\ref{ch:15} & メタアナリシス & & & & & 4 & & 3 & & \\
\ref{ch:16} & 欠測データ・多重補完 & & & & 5c & & & & 3 & \\
\ref{ch:17} & ポアソン率・ベイズ & & & & & & 4 & & & 2\\
\bottomrule
\end{tabular}
\end{center}

\paragraph{読み取れること}
\begin{itemize}
\item 生存時間（第9〜11章）は2016〜2023年の毎回出題され，2025年は競合リスクの理論として出た．
\item 分割表（第4章）は形を変えて繰り返し出る．特に「対応のある2 値」（McNemar型）は2018・2019・2024年，
      「層別・MH」は2017・2021・2022年に出ている．
\item 近年は臨床試験の方法論（多重性・欠測・混合モデル・ベースライン調整）が増えている．
\end{itemize}

% ------------------------------------------------------------
\chapter*{推奨学習ロードマップ}
\addcontentsline{toc}{chapter}{推奨学習ロードマップ}
\markboth{学習ロードマップ}{学習ロードマップ}

\begin{center}
\begin{tikzpicture}[node distance=5mm and 4mm,
  st/.style={draw=bkblue,rounded corners=2pt,fill=bkblue!5,text width=62mm,align=left,font=\small,inner sep=4pt},
  lb/.style={font=\small\bfseries,text=bkorange,text width=22mm,align=right}]
\node[st](s1){第1章 読み方，第2章 デザインと効果指標，第4章 分割表，第3章 連続・順位検定};
\node[lb,left=of s1]{段階1\\基礎体力};
\node[st,below=of s1](s2){第9章 生存時間の基礎，第10章 Cox・RMST・指数モデル，第11章 競合リスク};
\node[lb,left=of s2]{段階2\\最頻出};
\node[st,below=of s2](s3){第5章 診断・ROC，第6章 ロジスティック回帰，第7章 交絡・層別・PS};
\node[lb,left=of s3]{段階3\\2 値の応用};
\node[st,below=of s3](s4){第12章 検出力・症例数，第13章 同等性，第14章 中間解析・多重性};
\node[lb,left=of s4]{段階4\\臨床試験};
\node[st,below=of s4](s5){第8章 縦断・混合，第15章 メタアナリシス，第16章 欠測，第17章 ベイズ};
\node[lb,left=of s5]{段階5\\近年の論点};
\node[st,below=of s5,fill=bkorange!8,draw=bkorange](s6){付録A 手法選択表，付録B 公式集，付録E チェックリスト，過去問演習};
\node[lb,left=of s6]{直前期};
\foreach \a/\b in {s1/s2,s2/s3,s3/s4,s4/s5,s5/s6}\draw[-{Stealth},thick,bkblue](\a)--(\b);
\end{tikzpicture}
\end{center}

\begin{itemize}
\item 段階1で，効果指標（RD・RR・OR・HR）と「独立か対応ありか」の区別を固める．以降のすべての章の前提になる．
\item 段階2は最頻出分野である．KM推定値の表の作成，ログランク検定の手計算，打ち切りのある指数分布の尤度，
      Cox部分尤度の書き下しは，時間を計って何度も練習する．
\item 段階3・4は，式の導出よりも「どの手法をなぜ使うか」「結果をどう解釈するか」の記述が得点源になる．
\item 段階5は出題回数こそ少ないが，近年の出題が多く，一度理解すれば短時間で得点できる．
\item 直前期は付録Eのチェックリストで，自力で導出できない項目を潰す．
\end{itemize}

\tableofcontents

\mainmatter
\pagenumbering{arabic}
